\section*{Chapter \ref{ch:b2:weather-systems} --- Weather Systems}

\begin{solution}{exo:b2:weather-systems:1}
(a) $f = 2 \times 7.292 \times 10^{-5} \times \sin\ang{50} = \qty{1.117e-4}{s^{-1}}$. (b) With
$\bar T = \qty{278}{K}$, $\tan\alpha = 1.117 \times 10^{-4} \times 278 \times 15/(9.81 \times 6)
= 0.00791$, about 1 in 126. (c) $1/0.00791 = 126$: some \qty{126}{km} behind the front. (d) The
slope is inversely proportional to the temperature contrast: it halves, to 1 in 253.
\end{solution}

\begin{solution}{exo:b2:weather-systems:2}
(a) $6 \times 150 = \qty{900}{km}$ ahead. (b) $900/25 = \qty{36}{h}$, a day and a half. (c) First
the high cirrus, then cirrostratus, then altostratus; the cloud thickens and lowers, and rain
sets in from the low, thick cloud as the front comes near.
\end{solution}

\begin{solution}{exo:b2:weather-systems:3}
The symbols point the way the front moves, and into the warmer air ahead of a cold front or the
colder air ahead of a warm front. (a) A cold front: the cold air lies to its north-west, behind
it; it moves south-east. (b) A warm front: the cold air lies to its north, ahead of it; it moves
north. (c) An \omterm{def:b2:weather-systems:fronts}{occluded front}, joining the low to the triple point where the cold and warm
fronts start; cold air lies on both sides of it, the warm air lifted above; it moves the way its
symbols point. (d) A \omterm{def:b2:weather-systems:fronts}{stationary front}: the triangles point south into the warm air, so the cold
air lies to the north; it barely moves.
\end{solution}

\begin{solution}{exo:b2:weather-systems:4}
(a) $\beta = 2 \times 7.292 \times 10^{-5} \times \cos\ang{45}/(6.371 \times 10^6) =
\qty{1.62e-11}{m^{-1}.s^{-1}}$. (b) For \qty{4000}{km}, $k = \qty{1.571e-6}{m^{-1}}$, $\beta/k^2 =
\qty{6.6}{m/s}$ and $c = 12 - 6.6 = \qty{5.4}{m/s}$; for \qty{8000}{km}, $\beta/k^2 =
\qty{26.3}{m/s}$ and $c = \qty{-14.3}{m/s}$. (c) The short wave moves east, \qty{470}{km} a day;
the long one moves west, \qty{1230}{km} a day, against the wind.
\end{solution}

\begin{solution}{exo:b2:weather-systems:5}
(a) $c = 0$ when $U = \beta/k^2$, so $k = \sqrt{\beta/U}$ and $\lambda = 2\pi/k =
2\pi\sqrt{U/\beta}$. (b) $2\pi\sqrt{12/(1.62 \times 10^{-11})} = \qty{5410}{km}$. (c) The circle
of latitude \ang{45} measures $2\pi \times 6371 \times \cos\ang{45} = \qty{28300}{km}$: about 5.2
waves, so five in practice. (d) $\lambda$ grows as $\sqrt U$: a stronger westerly makes the
standing waves longer, and fewer of them fit round the hemisphere.
\end{solution}

\begin{solution}{exo:b2:weather-systems:6}
(a) At $C_+ = (\sqrt{72}, \sqrt{72}, 27)$: $\dot x = 10(\sqrt{72} - \sqrt{72}) = 0$, $\dot y =
\sqrt{72}(28 - 27) - \sqrt{72} = 0$, $\dot z = 72 - \tfrac83 \times 27 = 72 - 72 = 0$.
(b) $\mathrm{e}^{-13.67} = 1.16 \times 10^{-6}$: the volume shrinks almost a millionfold in one
time unit. (c) A volume can shrink while it is stretched in one direction and squeezed much more
strongly in the others: the distance along the stretched direction grows, and two nearby
trajectories separate along it, while the volume between them vanishes.
\end{solution}

\begin{solution}{exo:b2:weather-systems:7}
(a) $\ln 2/0.906 = 0.765$ time units. (b) $\ln(10/10^{-8})/0.906 = \ln(10^9)/0.906 = 22.9$ time
units. (c) $\ln 1000/0.906 = 7.6$ time units: a thousandfold improvement of the initial state
buys only a third of the time found in (b), since $\ln 10^3/\ln 10^9 = 1/3$.
\end{solution}

\begin{solution}{exo:b2:weather-systems:8}
(a) $(303 - 195)/195 = 0.554$ and $V_{\max} = \sqrt{0.9 \times 0.554 \times 15\,000} =
\qty{86.5}{m/s}$. (b) Well above the \qty{70}{m/s} of category 5, close to the \qty{86}{m/s} that
five recent storms have exceeded. (c) Storms rarely stay long in the steady state the theory
assumes; vertical wind shear, dry air drawn into the storm and the sinking air around it, which
the theory leaves out, weaken them.
\end{solution}

\begin{solution}{exo:b2:weather-systems:9}
(a) $(303 - 195)/303 = 0.356$. (b) $(303 - 195)/195 = 0.554$; the ratio is $0.554/0.356 =
303/195 = 1.554$. (c) In a Carnot engine the work leaves the engine. In a tropical cyclone the
work is dissipated by friction near the sea surface, and that heat warms the inflowing air at the
temperature $T_s$, where the engine takes in its heat: it is used again, and the engine behaves
as if it were more efficient by the factor $T_s/T_o$.
\end{solution}

\begin{solution}{exo:b2:weather-systems:10}
(a) The derivation of \cref{prop:b2:weather-systems:slope} without approximation gives
$\mathrm{d}z/\mathrm{d}x = f(\rho_c v_c - \rho_w v_w)/[g(\rho_c - \rho_w)]$. A denser fluid
resting on a lighter one is unstable: the slightest disturbance makes it overturn. In a front at
rest the cold, denser air must therefore lie beneath, which is the wedge of the definition; then
$\rho_c > \rho_w$, the surface rises towards the cold air, and $\rho_c v_c > \rho_w v_w$.
(b) $\zeta = \partial v/\partial x = 10/(50 \times 10^3) = \qty{2e-4}{s^{-1}}$, about $1.9f$ with
$f = \qty{1.03e-4}{s^{-1}}$. (c) Positive in the Northern Hemisphere: cyclonic, and as large as
the Coriolis parameter itself; fronts are bands of strong cyclonic vorticity.
\end{solution}

\begin{solution}{exo:b2:weather-systems:11}
(a) $\log_2(6/0.5) = \log_2 12 = 3.58$ doublings, $3.58 \times 2.4 = 8.6$~days. (b) One
more day of forecast needs one doubling time less of growth before the error is useless: the
initial error must shrink by $2^{1/2.4} = 1.33$ each decade. (c) A tenfold reduction is
$\log_2 10 = 3.32$ doublings; if it gains only 4 to 5 days, small errors double in $4/3.32 = 1.2$
to $5/3.32 = 1.5$ days: smaller errors grow faster than the larger ones measured in the
forecasts of the 1980s, through latent heat released in convection on small scales.
\end{solution}

\begin{solution}{exo:b2:weather-systems:12}
(a) At \qty{30}{\celsius}, $e_s = \qty{43.5}{hPa}$ and $q_s^* = \qty{27.1}{g/kg}$; the air at
\qty{29}{\celsius} has $e = 0.8 \times 41.0 = \qty{32.8}{hPa}$ and $q_a = \qty{20.4}{g/kg}$;
$h_o^* - h^* = 1004.5 \times 1 + 2.501 \times 10^6 \times 0.0067 = \qty{17.8}{kJ/kg}$, against
\qty{16.8}{kJ/kg}. (b) $(303.15 - 200)/200 = 0.516$ and $V_{\max} = \sqrt{0.9 \times 0.516
\times 17\,800} = \qty{90.9}{m/s}$, against \qty{87.9}{m/s}. (c) About \qty{3.4}{\percent} per
kelvin. The disequilibrium rises by \qty{6}{\percent}, which adds \qty{3}{\percent} to the speed;
the factor $(T_s - T_o)/T_o$ rises by \qty{1}{\percent}, which adds \qty{0.5}{\percent}: the
fuel, the moisture the warmer sea can give, contributes most, through the exponential rise of
$e_s$.
\end{solution}

\begin{solution}{pb:b2:weather-systems:1}
\textbf{1.} $(480 - 390)/6 = \qty{15}{km/h}$.
\textbf{2.} $390/15 = \qty{26}{h}$ after 12 UTC: at 14 UTC the next day.
\textbf{3.} $f = \qty{1.117e-4}{s^{-1}}$.
\textbf{4.} With $\bar T = \qty{281}{K}$, $\tan\alpha = 1.117 \times 10^{-4} \times 281 \times
16/(9.81 \times 8) = 0.0064$, about 1 in 156.
\textbf{5.} The front is then $6 \times 15 = \qty{90}{km}$ beyond the town; the surface is
$90\,000 \times 0.0064 = \qty{576}{m}$ above it.
\textbf{6.} $1500/0.0064 = \qty{234}{km}$ behind the front, reached $234/15 = \qty{15.6}{h}$ after
the passage.
\textbf{7.} Cumulus and cumulonimbus, showers, perhaps thunder, gusts and a sudden fall of
temperature; the pressure, falling as the front approaches, starts to rise as it passes; then
cooler, clearer air.
\textbf{8.} $\beta = 2 \times 7.292 \times 10^{-5} \times \cos\ang{50}/(6.371 \times 10^6) =
\qty{1.47e-11}{m^{-1}.s^{-1}}$.
\textbf{9.} $k = 2\pi/(5 \times 10^6) = \qty{1.257e-6}{m^{-1}}$ and $\beta/k^2 = \qty{9.3}{m/s}$.
\textbf{10.} $c = 20 - 9.3 = \qty{10.7}{m/s}$, eastward.
\textbf{11.} $10.7 \times 86\,400 = \qty{924}{km}$ east.
\textbf{12.} $\lambda_s = 2\pi\sqrt{20/(1.47 \times 10^{-11})} = \qty{7330}{km}$; the circle of
latitude \ang{50} measures $2\pi \times 6371 \times \cos\ang{50} = \qty{25700}{km}$, about 3.5
such waves.
\textbf{13.} Ahead of it, on its eastern side, where the short waves induce rising air and the
surface pressure falls.
\textbf{14.} $\lambda = \ln(3.9 \times 10^{-3}/5 \times 10^{-7})/(14 - 4) = \ln 7800/10 = 0.896$.
\textbf{15.} $\ln 2/0.896 = 0.77$ time units.
\textbf{16.} \qty{1}{\percent} below 0.906: the mean of 200 pairs, a finite time and readings from
a graph.
\textbf{17.} Once the two runs are far apart, they lie at unrelated points of the attractor, and
their distance cannot exceed its size: the growth saturates.
\textbf{18.} $\ln 10/0.896 = 2.6$ time units later.
\textbf{19.} $\log_2(3/0.1) = \log_2 30 = 4.91$ doublings.
\textbf{20.} $4.91 \times 2.4 = 11.8$~days.
\textbf{21.} One doubling time: 2.4~days.
\textbf{22.} $4.91 \times 5 = 24.5$~days: the early estimate was far too optimistic.
\textbf{23.} With a doubling time of 2.4 days, a tenfold smaller error would gain $3.32 \times 2.4
= 8.0$~days, not 4 to 5: small errors grow faster than 2.4 days per doubling, and an
estimate built on that doubling time overstates the horizon.
\textbf{24.} This forecast could be useful for at most about \textbf{12 days}, and in practice
for less.
\end{solution}
